% Abstract and keywords (the same as in the thesis)
\begin{abstract}
An unmanned surface vessel (USV) and an unmanned aerial vehicle (UAV) connected by a tether combine the endurance of the USV with the view from above of the UAV, which draws power and a reliable communication link through the tether. The tether, however, couples the two vehicles, which must remain within its reach while each respects its own limits, under waves and wind. This work develops the cooperative control of such a system. Models of both vehicles are derived and validated against a simulator, together with a closed-form model of the force of a slack tether. Each vehicle is driven by its own non-linear model predictive controller, and the two controllers exchange their predicted trajectories, so that each anticipates the motion of the other within the reach of the tether. In a simulation at sea, each controller tracks the trajectory given to its vehicle with an error of a few centimetres, also under wind, and, when the reference cannot be followed, both keep their vehicles within their limits and the tether within its reach. The UAV lands on the moving deck a few centimetres from the centre of the pad in waves of up to $1.2\,\text{m}$ and winds of up to $8\,\text{m/s}$, and the two vehicles operate as a single system to perform missions.
\end{abstract}

\begin{IEEEkeywords}
Cooperative control, non-linear model predictive control, tether, unmanned aerial vehicle, unmanned surface vessel
\end{IEEEkeywords}
